\FigureLayoutDeclare{img/2011/shaanxi_science/q16_fig1.png}{10.5cm}{6bp 55bp 72bp 26bp}

\chapter{2011年陕西卷（理）}

\section{单选题}

\begin{problem}
设 \(\bs{a}\)，\(\bs{b}\) 是向量，命题“若 \(\bs{a} = -\bs{b}\)，则 \(|\bs{a}| = |\bs{b}|\)”的逆命题是（\quad）

\choices
    {若 \(\bs{a} \neq -\bs{b}\)，则 \(|\bs{a}| \neq |\bs{b}|\).}
    {若 \(\bs{a} = -\bs{b}\)，则 \(|\bs{a}| \neq |\bs{b}|\)}
    {若 \(|\bs{a}| \neq |\bs{b}|\)，则 \(\bs{a} \neq -\bs{b}\)}
    {若 \(|\bs{a}| = |\bs{b}|\)，则 \(\bs{a} = -\bs{b}\)}
\end{problem}

\begin{answer}
D
\end{answer}

\begin{solution}
原命题的条件是“\(\bs{a}=-\bs{b}\)”，结论是“\(|\bs{a}|=|\bs{b}|\)”. 逆命题交换原命题的条件和结论，因此为“若 \(|\bs{a}|=|\bs{b}|\)，则 \(\bs{a}=-\bs{b}\)”，对应选项 D.
\end{solution}

\begin{problem}
设抛物线的顶点在原点，准线方程为 \(x = -2\)，则抛物线的方程是（\quad）

\choices
    {\( y^{2} = -8x \)}
    {\( y^{2} = -4x \)}
    {\( y^{2} = 8x \)}
    {\( y^{2} = 4x \)}
\end{problem}

\begin{answer}
C
\end{answer}

\begin{solution}
抛物线的顶点为原点，准线为 \(x=-2\)，所以焦点为 \(F(2,0)\)，抛物线开口向右. 设抛物线上任一点为 \(P(x,y)\)，由抛物线的定义，
\[
|PF|=|x+2|
\]
因此
\[
\sqrt{(x-2)^2+y^2}=|x+2|
\]
两边平方并化简，得 \(y^2=8x\)，故选 C.
\end{solution}

\begin{problem}
设函数 \( f(x) \) （\(x \in \R\)）满足 \(f(-x) = f(x)\)，\(f(x + 2) = f(x)\)，则 \( y = f(x) \) 的图象可能是（\quad）

\choices
    {\choicebitmap[width=0.15\paperwidth]{img/2011/shaanxi_science/q3_optA_clean.png}}
    {\choicebitmap[width=0.15\paperwidth]{img/2011/shaanxi_science/q3_optB_clean.png}}
    {\choicebitmap[width=0.15\paperwidth]{img/2011/shaanxi_science/q3_optC_clean.png}}
    {\choicebitmap[width=0.15\paperwidth]{img/2011/shaanxi_science/q3_optD_clean.png}}
\end{problem}

\begin{answer}
B
\end{answer}

\begin{solution}
由 \(f(-x)=f(x)\)，函数图象关于 \(y\) 轴对称；由 \(f(x+2)=f(x)\)，函数图象向右平移 \(2\) 个单位后与原图象重合. 选项 B 的图象同时具有这两个性质，例如其图象可由 \(y=-\cos \pi x\) 表示. 其余图象至少有一个性质不满足，因此选 B.
\end{solution}

\begin{problem}
\((4^{x} - 2^{-x})^{6}\)（\(x \in \R\)）展开式中的常数项是（\quad）

\choices
    {\(-20\)}
    {\(-15\)}
    {15}
    {20}
\end{problem}

\begin{answer}
C
\end{answer}

\begin{solution}
令 \(t=2^x\)，则 \(t>0\)，且 \(4^x=t^2\)，\(2^{-x}=t^{-1}\). 展开式的一般项为
\[
\mathrm C_6^k(t^2)^{6-k}(-t^{-1})^k=(-1)^k\mathrm C_6^k t^{12-3k},\quad k=0,1,\ldots,6
\]
要得到常数项，必须有 \(12-3k=0\)，所以 \(k=4\). 常数项为
\[
(-1)^4\mathrm C_6^4=15
\]
故选 C.
\end{solution}

\begin{problem}
某几何体的三视图如图所示，则它的体积是（\quad）

\bitmapfigure[max width=0.52\linewidth]{img/2011/shaanxi_science/q5_fig1.png}


\choices
    {\(8 - \frac{2\pi}{3}\)}
    {\(8 - \frac{\pi}{3}\)}
    {\(8 - 2\pi\)}
    {\(\frac{2\pi}{3}\)}
\end{problem}

\begin{answer}
A
\end{answer}

\begin{solution}
由三视图可知，该几何体可看作棱长为 \(2\) 的正方体内部去掉一个底面半径为 \(1\)、高为 \(2\) 的圆锥. 因此正方体体积为 \(2^3=8\)，圆锥体积为
\[
\frac{1}{3}\pi\cdot 1^2\cdot 2=\frac{2\pi}{3}
\]
所以几何体体积为
\[
8-\frac{2\pi}{3}
\]
故选 A.
\end{solution}

\begin{problem}
函数 \( f(x) = \sqrt{x} - \cos x \) 在 \([0, +\infty)\) 内（\quad）

\choices
    {没有零点}
    {有且仅有一个零点}
    {有且仅有两个零点}
    {有无穷多个零点}
\end{problem}

\begin{answer}
B
\end{answer}

\begin{solution}
函数在 \([0,+\infty)\) 上连续，且 \(f(0)=-1<0\)，\(f(1)=1-\cos 1>0\).
由零点存在定理，函数在 \((0,1)\) 内至少有一个零点. 对于 \(0<x<1\)，
\[
f'(x)=\frac{1}{2\sqrt{x}}+\sin x>0
\]
所以 \(f(x)\) 在 \((0,1)\) 上严格递增，因而在该区间内至多有一个零点.

当 \(x>1\) 时，\(\sqrt{x}>1\geqslant\cos x\)，从而 \(f(x)>0\)；而 \(f(1)>0\). 所以 \([1,+\infty)\) 内没有零点. 因此函数在 \([0,+\infty)\) 内有且仅有一个零点，故选 B.
\end{solution}

\begin{problem}
设集合 \(M = \{y \mid y = |\cos^2 x - \sin^2 x|, x \in \R\}\)，\(N = \left\{x \mid \left|x - \frac{1}{\i}\right| < \sqrt{2}, x \in \R\right\}\)，其中 \(\i\) 为虚数单位，则 \(M \cap N\) 为（\quad）

\choices
    {\((0,1)\)}
    {\((0,1]\)}
    {\([0,1)\)}
    {\([0,1]\)}
\end{problem}

\begin{answer}
C
\end{answer}

\begin{solution}
因为 \(\frac{1}{\i}=-\i\)，对实数 \(x\) 有
\[
\left|x-\frac{1}{\i}\right|=|x+\i|=\sqrt{x^2+1}
\]
于是
\[
\sqrt{x^2+1}<\sqrt{2}\Longleftrightarrow x^2<1\Longleftrightarrow -1<x<1
\]
所以 \(N=(-1,1)\).

又因为
\[
|\cos^2x-\sin^2x|=|\cos 2x|
\]
且 \(|\cos 2x|\) 的值域为 \([0,1]\)，所以 \(M=[0,1]\). 因此
\[
M\cap N=[0,1)
\]
故选 C.
\end{solution}

\begin{problem}
如图，\(x_{1}\)，\(x_{2}\)，\(x_{3}\) 为某次考试三个评阅人对同一道题的独立评分，\(p\) 为该题的最终得分，当 \(x_{1} = 6\)，\(x_{2} = 9\)，\(p = 8.5\) 时，\(x_{3}\) 等于（\quad）

\bitmapfigure[max width=0.34\linewidth]{img/2011/shaanxi_science/q8_fig1.png}

\choices
    {11}
    {10}
    {8}
    {7}
\end{problem}

\begin{answer}
C
\end{answer}

\begin{solution}
由流程图，先计算 \(|x_1-x_2|=|6-9|=3>2\)，所以必须输入 \(x_3\)，不会直接输出 \(p\).

若满足 \(|x_3-x_1|<|x_3-x_2|\)，流程图令 \(x_2=x_3\)，于是
\[
\frac{x_1+x_3}{2}=8.5\Longrightarrow x_3=11
\]
但此时 \(|x_3- x_1|=5>|x_3-x_2|=2\)，不满足该分支条件，矛盾.

因此应进入另一分支，令 \(x_1=x_3\). 此时
\[
\frac{x_3+x_2}{2}=8.5\Longrightarrow \frac{x_3+9}{2}=8.5\Longrightarrow x_3=8
\]
检验得 \(|8-6|=2\geqslant|8-9|=1\)，确实进入该分支，故选 C.
\end{solution}

\begin{problem}
设 \((x_{1}, y_{1})\)，\((x_{2}, y_{2})\)，\(\cdots\)，\((x_{n}, y_{n})\) 是变量 \(x\) 和 \(y\) 的 \(n\) 个样本，直线 \(l\) 是由这些样本点通过最小二乘法得到的线性回归直线（如图），以下结论中正确的是（\quad）

\bitmapfigure[max width=0.44\linewidth]{img/2011/shaanxi_science/q9_fig1.png}

\choices
    {\(x\) 和 \(y\) 的相关系数为直线 \(l\) 的斜率}
    {\(x\) 和 \(y\) 的相关系数在 \(0\) 到 \(1\) 之间}
    {当 \(n\) 为偶数时，分布在 \(l\) 两侧的样本点的个数一定相同}
    {直线 \(l\) 过点 \((\bar{x}, \bar{y})\)}
\end{problem}

\begin{answer}
D
\end{answer}

\begin{solution}
设线性回归直线为 \(y=\hat{a}x+\hat{b}\). 最小二乘法使
\[
\sum_{i=1}^{n}(y_i-\hat{a}x_i-\hat{b})^2
\]
最小. 对 \(\hat{b}\) 求偏导并令其为零，得
\[
\sum_{i=1}^{n}(y_i-\hat{a}x_i-\hat{b})=0
\]
两边除以 \(n\)，得到 \(\bar{y}=\hat{a}\bar{x}+\hat{b}\)，所以回归直线经过点 \((\bar{x},\bar{y})\).

相关系数与回归直线斜率一般不相等；相关系数的取值范围是 \([-1,1]\)，不一定在 \([0,1]\) 内；样本点在回归直线两侧的个数也不由 \(n\) 的奇偶性决定. 因此正确选项为 D.
\end{solution}

\begin{problem}
甲、乙两人一起去游“2011西安世园会”，他们约定：各自独立地从 \(1\) 到 \(6\) 号景点中任选 \(4\) 个进行游览，每个景点参观 \(1\text{ 小时}\)，则最后 \(1\text{ 小时}\) 他们同在一个景点的概率是（\quad）

\choices
    {\( \frac{1}{36} \)}
    {\( \frac{1}{9} \)}
    {\( \frac{5}{36} \)}
    {\( \frac{1}{6} \)}
\end{problem}

\begin{answer}
D
\end{answer}

\begin{solution}
每个人最后一小时所游览的景点在 \(1\) 至 \(6\) 号景点中具有相同的可能性，因此甲、乙最后一小时所在景点分别独立且均匀地取自这 \(6\) 个景点. 固定甲最后一小时所在的景点，乙与甲相同的概率为 \(\frac{1}{6}\). 因此
\[
P(\text{两人最后一小时同在一个景点})=\frac{1}{6}
\]
故选 D.
\end{solution}

\section{填空题}

\begin{problem}
设 \(f(x)=\begin{cases}
\lg x, & x>0,\\
x+\int_{0}^{a} 3t^{2}\,dt, & x\leqslant 0
\end{cases}\)，若 \(f(f(1))=1\)，则 \(a=\fillinblank{}\).
\end{problem}

\begin{answer}
1
\end{answer}

\begin{solution}
因为 \(1>0\)，所以 \(f(1)=\lg 1=0\). 再由第二段定义，
\[
f(f(1))=f(0)=0+\int_{0}^{a}3t^2\,dt=a^3
\]
由题设 \(f(f(1))=1\)，得 \(a^3=1\). 因为 \(a\) 为实数，所以 \(a=1\).
\end{solution}

\begin{problem}
设 \(n \in \N^*\)，一元二次方程 \(x^2 - 4x + n = 0\) 有整数根的充要条件是 \(n = \)\fillinblank{}.
\end{problem}

\begin{answer}
\(3\) 或 \(4\)
\end{answer}

\begin{solution}
设方程的两个整数根为 \(r\)，\(s\). 由韦达定理，
\[
\begin{cases}
r+s=4,\\
rs=n
\end{cases}
\]
因为 \(n\in\N^*\)，所以 \(rs>0\)，再结合 \(r+s=4\)，可知 \(r\)，\(s\) 均为正整数. 正整数和为 \(4\) 时只有
\[
\{r,s\}=\{1,3\}\quad\text{或}\quad \{r,s\}=\{2,2\}
\]
于是 \(n=3\) 或 \(n=4\). 反之，\(n=3\) 时方程为 \((x-1)(x-3)=0\)，\(n=4\) 时方程为 \((x-2)^2=0\)，均有整数根. 因此充要条件是 \(n=3\) 或 \(n=4\).
\end{solution}

\begin{problem}
观察下列等式 \(1 = 1\)，\(2 + 3 + 4 = 9\)，\(3 + 4 + 5 + 6 + 7 = 25\)，\(4 + 5 + 6 + 7 + 8 + 9 + 10 = 49\). 照此规律，第 \(n\) 个等式应为\fillinblank{}.
\end{problem}

\begin{answer}
\(n+(n+1)+(n+2)+\cdots+(3n-2)=(2n-1)^2\)
\end{answer}

\begin{solution}
第 \(n\) 个等式左边从 \(n\) 开始，每次增加 \(1\)，并且项数依次为 \(1,3,5,7,\ldots\)，所以第 \(n\) 个等式有 \(2n-1\) 项，最后一项为
\[
n+(2n-2)=3n-2
\]
因此左边是首项为 \(n\)、末项为 \(3n-2\)、项数为 \(2n-1\) 的等差数列，其和为
\[
\frac{\bigl[n+(3n-2)\bigr](2n-1)}{2}=(2n-1)^2
\]
所以第 \(n\) 个等式为
\[
n+(n+1)+(n+2)+\cdots+(3n-2)=(2n-1)^2
\]
\end{solution}

\begin{problem}
植树节某班 \(20\) 名同学在一段直线公路一侧植树，每人植一棵，相邻两棵树相距 \(10\text{ 米}\)，开始时需将树苗集中放置在某一树坑旁边，使每位同学从各自树坑出发前来领取树苗往返所走的路程总和最小，这个最小值为\fillinblank{}\(\text{米}\).
\end{problem}

\begin{answer}
\(2000\)
\end{answer}

\begin{solution}
将 \(20\) 个树坑依次编号为 \(1,2,\ldots,20\)，相邻树坑距离为 \(10\text{ 米}\). 若树苗集中放在第 \(k\) 个树坑旁，则所有同学领取树苗后往返所走的总路程为
\[
S_k=2\cdot10\sum_{i=1}^{20}|i-k|=20\sum_{i=1}^{20}|i-k|
\]
令 \(T_k=\sum_{i=1}^{20}|i-k|\). 当 \(1\leqslant k\leqslant19\) 时，树苗位置从第 \(k\) 个树坑移到第 \(k+1\) 个树坑，左侧 \(k\) 棵树的距离总和增加 \(k\)，右侧 \(20-k\) 棵树的距离总和减少 \(20-k\)，所以
\[
T_{k+1}-T_k=k-(20-k)=2k-20
\]
因此 \(T_k\) 在 \(k=10\) 与 \(k=11\) 处取得最小值. 取 \(k=10\)，有
\[
T_{10}=(1+2+\cdots+9)+(1+2+\cdots+10)=45+55=100
\]
故最小总路程为
\[
S_{10}=20\times100=2000\text{ 米}
\]
\end{solution}

\section{选做题}

\begin{problem}
若关于 \(x\) 的不等式 \(\left|a\right| \geqslant \left|x + 1\right| + \left|x - 2\right|\) 存在实数解，则实数 \(a\) 的取值范围是\fillinblank{}.
\end{problem}

\begin{answer}
\((-\infty,-3]\cup[3,+\infty)\)
\end{answer}

\begin{solution}
由三角不等式，任意实数 \(x\) 都有
\[
|x+1|+|x-2|\geqslant|(x+1)-(x-2)|=3
\]
当且仅当 \(-1\leqslant x\leqslant2\) 时取等号，因此不等式存在实数解的充要条件是 \(|a|\geqslant3\). 所以
\[
a\in(-\infty,-3]\cup[3,+\infty)
\]
\end{solution}

\begin{problem}
如图，\(\angle B = \angle D\)，\(AE \perp BC\)，\(\angle ACD = 90^{\circ}\)，且 \(AB = 6\)，\(AC = 4\)，\(AD = 12\)，则 \(BE =\fillinblank{}\).

\bitmapfigure[max width=0.30\linewidth]{img/2011/shaanxi_science/q15_fig1.png}
\end{problem}

\begin{answer}
\(4\sqrt{2}\)
\end{answer}

\begin{solution}
由图，\(E\) 是 \(AD\) 与 \(BC\) 的交点. 在直角三角形 \(ACD\) 中，\(\angle ACD=90^\circ\)，且 \(AD=12\)，\(AC=4\)，所以
\[
DC=\sqrt{AD^2-AC^2}=\sqrt{12^2-4^2}=8\sqrt{2}
\]
因此
\[
\cos\angle D=\cos\angle ADC=\frac{DC}{AD}=\frac{2\sqrt{2}}{3}
\]
又因为 \(\angle B=\angle D\)，且 \(BE\subset BC\)，所以 \(\angle ABE=\angle B=\angle D\). 由于 \(AE\perp BC\)，三角形 \(ABE\) 在 \(E\) 处为直角，\(AB=6\) 为斜边，于是
\[
BE=AB\cos\angle ABE=6\cdot\frac{2\sqrt{2}}{3}=4\sqrt{2}
\]
\end{solution}

\begin{problem}
直角坐标系 \(xOy\) 中，以原点为极点，\(x\) 轴正半轴为极轴建立极坐标系，设点 \(A\)，\(B\) 分别在曲线 \( C_{1}:\begin{cases}
x=3+\cos\theta,\\
y=4+\sin\theta,
\end{cases} \)（\(\theta\) 为参数）和曲线 \( C_{2}:\rho=1 \) 上，则 \( |AB| \) 的最小值为\fillinblank{}.
\end{problem}

\begin{answer}
3
\end{answer}

\begin{solution}
参数方程
\[
\begin{cases}
x=3+\cos\theta,\\
y=4+\sin\theta
\end{cases}
\]
表示圆 \(C_1\)，其圆心为 \(O_1(3,4)\)，半径为 \(1\). 极坐标方程 \(\rho=1\) 表示以原点 \(O\) 为圆心、半径为 \(1\) 的圆 \(C_2\). 两圆心距离为
\[
OO_1=\sqrt{3^2+4^2}=5
\]
对任意 \(A\in C_1\)，\(B\in C_2\)，由三角不等式，
\[
5=|OO_1|\leqslant|OB|+|BA|+|AO_1|=1+|AB|+1
\]
所以 \(|AB|\geqslant3\). 当 \(O\)，\(B\)，\(A\)，\(O_1\) 依次在两圆心连线上且 \(B\)，\(A\) 分别为两圆相向的内侧点时取等号，故 \(|AB|\) 的最小值为 \(3\).
\end{solution}

\section{解答题}

\begin{problem}
如图，在 \(\triangle ABC\) 中，\(\angle ABC = 60^{\circ}\)，\(\angle BAC = 90^{\circ}\)，\(AD\) 是 \(BC\) 上的高，沿 \(AD\) 把 \(\triangle ABD\) 折起，使 \(\angle BDC = 90^{\circ}\).
\begin{enumerate}
\item 证明：平面 \(ADB\perp\) 平面 \(BDC\)；
\item 设 \(E\) 为 \(BC\) 的中点，求 \(\overrightarrow{AE}\) 与 \(\overrightarrow{DB}\) 夹角的余弦值.
\end{enumerate}

\bitmapfigure[max width=0.32\linewidth]{img/2011/shaanxi_science/q16_fig1.png}
\end{problem}

\begin{solution}
\begin{enumerate}
\item 因为 \(AD\) 是 \(BC\) 上的高，所以 \(AD\perp BC\)，从而 \(AD\perp DC\). 又因为折起后 \(\angle BDC=90^{\circ}\)，所以 \(BD\perp DC\). 直线 \(AD\)，\(BD\) 是平面 \(ADB\) 内相交的两条直线，且 \(DC\) 同时垂直于它们，因此 \(DC\perp\) 平面 \(ADB\). 又 \(DC\subset\) 平面 \(BDC\)，所以平面 \(ADB\perp\) 平面 \(BDC\).
\item 在折叠前的直角三角形 \(ABD\) 中，\(AD\perp BD\)，折叠后这一垂直关系保持不变. 在直角三角形 \(ABD\) 中，\(\angle ABD=60^{\circ}\)，故 \(\angle BAD=30^{\circ}\)，于是 \(AD=\sqrt{3}BD\). 由直角三角形斜边上的高的性质，\(AD^{2}=BD\cdot DC\)，所以 \(DC=3BD\). 不妨设 \(BD=1\)，建立直角坐标系，使 \(D\) 为原点，\(DC\) 在 \(x\) 轴正半轴上，\(DB\) 在 \(y\) 轴正半轴上，\(DA\) 在 \(z\) 轴正半轴上，则
\(D(0,0,0)\)，\(B(0,1,0)\)，\(C(3,0,0)\)，\(A(0,0,\sqrt{3})\).
设 \(E\) 为 \(BC\) 的中点，则 \(E\left(\frac{3}{2},\frac{1}{2},0\right)\)，从而
\(\overrightarrow{AE}=\left(\frac{3}{2},\frac{1}{2},-\sqrt{3}\right)\)，\(\overrightarrow{DB}=(0,1,0)\).
因此 \(\overrightarrow{AE}\cdot\overrightarrow{DB}=\frac{1}{2}\)，\(\left|\overrightarrow{DB}\right|=1\)，且 \(\left|\overrightarrow{AE}\right|^{2}=\frac{9}{4}+\frac{1}{4}+3=\frac{11}{2}\). 设所求夹角为 \(\theta\)，则
\[
\cos\theta=\frac{\overrightarrow{AE}\cdot\overrightarrow{DB}}{\left|\overrightarrow{AE}\right|\left|\overrightarrow{DB}\right|}=\frac{\frac{1}{2}}{\sqrt{\frac{11}{2}}}=\frac{1}{\sqrt{22}}
\]
\end{enumerate}
\end{solution}

\begin{problem}
如图，设 \(P\) 是圆 \(x^{2} + y^{2} = 25\) 上的动点，点 \(D\) 是 \(P\) 在 \(x\) 轴上的投影，\(M\) 为 \(PD\) 上一点，且 \(|MD| = \frac{4}{5} |PD|\).
\begin{enumerate}
\item 当 \(P\) 在圆上运动时，求点 \(M\) 的轨迹 \(C\) 的方程；
\item 求过点 \( (3,0) \) 且斜率为 \( \frac{4}{5} \) 的直线被 \(C\) 所截线段的长度.
\end{enumerate}

\bitmapfigure[max width=0.36\linewidth]{img/2011/shaanxi_science/q17_fig1.png}
\end{problem}

\begin{solution}
\begin{enumerate}
\item 设 \(P(u,v)\)，则 \(D(u,0)\). 由于 \(M\) 在线段 \(PD\) 上且 \(|MD|=\frac{4}{5}|PD|\)，有 \(M\left(u,\frac{4}{5}v\right)\). 设 \(M(x,y)\)，则 \(u=x\)，\(v=\frac{5}{4}y\). 代入圆 \(P\) 的方程，得
\[
x^{2}+\frac{25}{16}y^{2}=25
\]
即
\[
C:\quad \frac{x^{2}}{25}+\frac{y^{2}}{16}=1
\]
反之，曲线 \(C\) 上任一点 \((x,y)\) 取 \(P\left(x,\frac{5}{4}y\right)\)，便满足 \(x^{2}+\frac{25}{16}y^{2}=25\)，因此该点确由某个圆上的 \(P\) 产生，故这就是点 \(M\) 的完整轨迹.
\item 过点 \((3,0)\) 且斜率为 \(\frac{4}{5}\) 的直线为 \(y=\frac{4}{5}(x-3)\). 代入 \(C\) 的方程，得
\[
\frac{x^{2}}{25}+\frac{1}{16}\left[\frac{4}{5}(x-3)\right]^{2}=1
\]
即 \(x^{2}+(x-3)^{2}=25\)，所以 \(x^{2}-3x-8=0\)，两交点的横坐标之差为 \(\sqrt{41}\). 由于直线斜率为 \(\frac{4}{5}\)，两交点的纵坐标之差为 \(\frac{4}{5}\sqrt{41}\)，故所截线段长度为
\[
\sqrt{(\sqrt{41})^{2}+\left(\frac{4}{5}\sqrt{41}\right)^{2}}=\frac{41}{5}
\]
\end{enumerate}
\end{solution}

\begin{problem}
叙述并证明余弦定理.
\end{problem}

\begin{solution}
在 \(\triangle ABC\) 中，设 \(a=|BC|\)，\(b=|CA|\)，\(c=|AB|\). 建立直角坐标系，令 \(A\) 为原点，\(AB\) 在 \(x\) 轴正半轴上，则可取 \(B(c,0)\)，\(C(b\cos A,b\sin A)\). 因此
\[
a^{2}=|BC|^{2}=(b\cos A-c)^{2}+(b\sin A)^{2}=b^{2}\cos^{2}A-2bc\cos A+c^{2}+b^{2}\sin^{2}A=b^{2}+c^{2}-2bc\cos A
\]

再以 \(B\) 为原点，令 \(BC\) 在 \(x\) 轴正半轴上，则可取 \(C(a,0)\)，\(A(c\cos B,c\sin B)\). 因此
\[
b^{2}=|CA|^{2}=(c\cos B-a)^{2}+(c\sin B)^{2}=c^{2}\cos^{2}B-2ac\cos B+a^{2}+c^{2}\sin^{2}B=c^{2}+a^{2}-2ac\cos B
\]

最后以 \(C\) 为原点，令 \(CA\) 在 \(x\) 轴正半轴上，则可取 \(A(b,0)\)，\(B(a\cos C,a\sin C)\). 因此
\[
c^{2}=|AB|^{2}=(a\cos C-b)^{2}+(a\sin C)^{2}=a^{2}\cos^{2}C-2ab\cos C+b^{2}+a^{2}\sin^{2}C=a^{2}+b^{2}-2ab\cos C
\]
以上三个等式即为余弦定理.
\end{solution}

\begin{problem}
如图，从点 \(P_{1}(0,0)\) 作 \(x\) 轴的垂线交曲线 \(y = \e^{x}\) 于点 \(Q_{1}(0,1)\)，曲线在 \(Q_{1}\) 点处的切线与 \(x\) 轴交于点 \(P_{2}\)，再从 \(P_{2}\) 作 \(x\) 轴的垂线交曲线于点 \(Q_{2}\)，依次重复上述过程得到一系列点：\(P_{1}\)，\(Q_{1}\)；\(P_{2}\)，\(Q_{2}\)；\(\cdots\)；\(P_{n}\)，\(Q_{n}\)，记 \(P_{k}\) 点的坐标为 \((x_{k}, 0)\)（\(k = 1, 2, \cdots, n\)）.
\begin{enumerate}
\item 试求 \(x_{k}\) 与 \(x_{k-1}\) 的关系（\(2 \leqslant k \leqslant n\)）；
\item 求 \(\left|P_{1}Q_{1}\right|+\left|P_{2}Q_{2}\right|+\left|P_{3}Q_{3}\right|+\cdots+\left|P_{n}Q_{n}\right|\).
\end{enumerate}

\bitmapfigure[max width=0.54\linewidth]{img/2011/shaanxi_science/q19_fig1.png}
\end{problem}

\begin{solution}
\begin{enumerate}
\item 对 \(Q_{k-1}\left(x_{k-1},\e^{x_{k-1}}\right)\) 处的曲线 \(y=\e^{x}\)，其切线方程为
\[
y-\e^{x_{k-1}}=\e^{x_{k-1}}(x-x_{k-1})
\]
令 \(y=0\)，并利用切线与 \(x\) 轴交于 \(P_k(x_k,0)\)，得
\[
-\e^{x_{k-1}}=\e^{x_{k-1}}(x_k-x_{k-1})
\]
所以 \(x_k=x_{k-1}-1\)（\(2\leqslant k\leqslant n\)）. 又 \(x_1=0\)，故 \(x_k=1-k\).
\item 因为 \(Q_k\left(x_k,\e^{x_k}\right)\)，且 \(P_k\) 在 \(x\) 轴上，所以 \(|P_kQ_k|=\e^{x_k}=\e^{1-k}\). 因此
\[
|P_1Q_1|+|P_2Q_2|+\cdots+|P_nQ_n|=1+\frac{1}{\e}+\frac{1}{\e^{2}}+\cdots+\frac{1}{\e^{n-1}}=\frac{1-\e^{-n}}{1-\e^{-1}}=\frac{\e(1-\e^{-n})}{\e-1}
\]
\end{enumerate}
\end{solution}

\begin{problem}
如图，\(A\) 地到火车站共有两条路径 \(L_{1}\) 和 \(L_{2}\)，据统计，通过两条路径所用的时间互不影响，所用时间落在各时间段内的频率如下表：
\begin{center}
\begin{tabular}{c|ccccc}
\hline
所用时间（分钟） & \(10\sim20\) & \(20\sim30\) & \(30\sim40\) & \(40\sim50\) & \(50\sim60\) \\
\hline
选择 \(L_1\) 的人数 & \(6\) & \(12\) & \(18\) & \(12\) & \(12\) \\
选择 \(L_2\) 的人数 & \(0\) & \(4\) & \(16\) & \(16\) & \(4\) \\
\hline
\end{tabular}
\end{center}
现甲、乙两人分别有 \(40\text{ 分钟}\) 和 \(50\text{ 分钟}\) 时间用于赶往火车站.
\begin{enumerate}
\item 为了尽最大可能在各自允许的时间内赶到火车站，甲和乙应如何选择各自的路径？
\item 用 \(X\) 表示甲、乙两人中在允许的时间内能赶到火车站的人数，针对（1）的选择方案，求 \(X\) 的分布列和数学期望.
\end{enumerate}

\begingroup
\leftskip=0pt\relax
\bitmapfigure[width=0.30\linewidth]{img/2011/shaanxi_science/q20_fig1.png}
\endgroup
\end{problem}

\begin{solution}
\begin{enumerate}
\item 选择 \(L_1\) 的总人数为 \(60\)，其中用时不超过 \(40\) 分钟的有 \(6+12+18=36\) 人，所以甲选择 \(L_1\) 时在允许时间内到达的概率为 \(\frac{36}{60}=\frac{3}{5}\). 选择 \(L_2\) 的总人数为 \(40\)，其中用时不超过 \(40\) 分钟的有 \(0+4+16=20\) 人，所以甲选择 \(L_2\) 时相应概率为 \(\frac{20}{40}=\frac{1}{2}\). 故甲应选择 \(L_1\).

选择 \(L_1\) 时用时不超过 \(50\) 分钟的有 \(6+12+18+12=48\) 人，概率为 \(\frac{48}{60}=\frac{4}{5}\)；选择 \(L_2\) 时用时不超过 \(50\) 分钟的有 \(0+4+16+16=36\) 人，概率为 \(\frac{36}{40}=\frac{9}{10}\). 故乙应选择 \(L_2\).
\item 设事件 \(A\) 为甲在允许时间内到达，事件 \(B\) 为乙在允许时间内到达. 由题意，\(A\)，\(B\) 相互独立，且 \(P(A)=\frac{3}{5}\)，\(P(B)=\frac{9}{10}\). 因为 \(X\) 表示两人中在允许时间内到达的人数，所以
\[
P(X=0)=P(\overline{A})P(\overline{B})=\frac{2}{5}\cdot\frac{1}{10}=\frac{1}{25}
\]
\[
P(X=1)=P(A)P(\overline{B})+P(\overline{A})P(B)=\frac{3}{5}\cdot\frac{1}{10}+\frac{2}{5}\cdot\frac{9}{10}=\frac{21}{50}
\]
\[
P(X=2)=P(A)P(B)=\frac{3}{5}\cdot\frac{9}{10}=\frac{27}{50}
\]
因此 \(X\) 的分布列为
\begin{center}
\begin{tabular}{c|ccc}
\hline
\(X\) & \(0\) & \(1\) & \(2\) \\
\hline
\(P\) & \(\frac{1}{25}\) & \(\frac{21}{50}\) & \(\frac{27}{50}\) \\
\hline
\end{tabular}
\end{center}
数学期望为
\[
E(X)=0\cdot\frac{1}{25}+1\cdot\frac{21}{50}+2\cdot\frac{27}{50}=\frac{75}{50}=\frac{3}{2}
\]
\end{enumerate}
\end{solution}

\begin{problem}
设函数 \(f(x)\) 定义在 \((0, +\infty)\) 上，\(f(1) = 0\)，导函数 \(f'(x) = \frac{1}{x}\)，\(g(x) = f(x) + f'(x)\).
\begin{enumerate}
\item 求 \(g(x)\) 的单调区间和最小值；
\item 讨论 \(g(x)\) 与 \(g\left(\frac{1}{x}\right)\) 的大小关系；
\item 是否存在 \(x_0 > 0\)，使得 \(|g(x) - g(x_0)| < \frac{1}{2}\) 对任意 \(x > 0\) 成立？若存在，求出 \(x_0\) 的取值范围；若不存在，请说明理由.
\end{enumerate}
\end{problem}

\begin{solution}
由 \(f'(x)=\frac{1}{x}\) 和 \(f(1)=0\)，得
\[
f(x)=\int_{1}^{x}\frac{1}{t}\,dt=\ln x
\]
所以 \(g(x)=\ln x+\frac{1}{x}\).

\begin{enumerate}
\item \(g'(x)=\frac{1}{x}-\frac{1}{x^{2}}=\frac{x-1}{x^{2}}\). 因为 \(x>0\)，当 \(0<x<1\) 时 \(g'(x)<0\)，当 \(x>1\) 时 \(g'(x)>0\)，故 \(g(x)\) 在 \((0,1)\) 上单调递减，在 \((1,+\infty)\) 上单调递增，最小值为 \(g(1)=1\).
\item 令 \(h(x)=g(x)-g\left(\frac{1}{x}\right)\)，则
\[
h(x)=2\ln x+\frac{1}{x}-x
\]
且
\[
h'(x)=\frac{2}{x}-\frac{1}{x^{2}}-1=-\frac{(x-1)^{2}}{x^{2}}\leqslant0
\]
因此 \(h(x)\) 在 \((0,+\infty)\) 上单调递减，且 \(h(1)=0\). 所以当 \(0<x<1\) 时 \(h(x)>0\)，当 \(x=1\) 时 \(h(x)=0\)，当 \(x>1\) 时 \(h(x)<0\). 即当 \(0<x<1\) 时 \(g(x)>g\left(\frac{1}{x}\right)\)，当 \(x=1\) 时 \(g(x)=g\left(\frac{1}{x}\right)\)，当 \(x>1\) 时 \(g(x)<g\left(\frac{1}{x}\right)\).
\item 由 \(g(x)=\ln x+\frac{1}{x}\)，有 \(\lim_{x\to0^{+}}g(x)=+\infty\). 对任意固定的 \(x_0>0\)，总可取充分接近 \(0\) 的 \(x>0\)，使 \(g(x)>g(x_0)+\frac{1}{2}\)，从而 \(\left|g(x)-g(x_0)\right|>\frac{1}{2}\). 因此不存在满足条件的 \(x_0\).
\end{enumerate}
\end{solution}
